\section{Phase Change: de la loss bump a la circuit hypothesis}

El hallazgo principal del loss map bidimensional no es que el ICL se fortalezca poco a poco, sino que aparece de repente dentro de una ventana de entrenamiento muy estrecha. Para resaltar este cambio, los investigadores calculan diferencias o derivadas de la token-position curve y de la training trajectory. En el mapa de calor aparece una línea vertical nítida, y la loss curve ordinaria también muestra una pequeña bump.

\lecturefigure{slide-08-per-step-learning-heatmaps.jpg}{Loss surface y vista de derivadas: dentro de una ventana de entrenamiento estrecha aparece un cambio estructural evidente.}{Video oficial de Stanford Online, 00:13:48--00:14:01.}

Las magnitudes de diagnóstico pueden escribirse de forma abstracta como
\[
D_i(s,i)=\frac{\partial \mathcal{L}(s,i)}{\partial \log i},
\qquad
D_s(s)=\frac{\partial S_{\mathrm{ICL}}(s)}{\partial \log s}.
\]

\begin{itemize}
  \item $D_i$ mide la loss reduction que produce aumentar la longitud relativa del context;
  \item $D_s$ mide la velocidad con la que cambia el ICL score durante el entrenamiento;
  \item un narrow spike indica que el comportamiento cambia rápidamente en pocos training steps;
  \item la derivada amplifica el ruido, por lo que es indispensable contrastarla con la loss original, con varias seeds y con varios model sizes.
\end{itemize}

\lecturefigure{slide-09-phase-change-across-scales.jpg}{En distintas model scales se observan la loss bump y el ascenso rápido del ICL.}{Video oficial de Stanford Online, 00:14:02--00:15:11.}

\begin{warningbox}{«Phase change» es una descripción empírica, no una transición de fase física demostrada}
Que la curva cambie rápidamente en un número finito de snapshots no demuestra que la optimization trajectory sea matemáticamente discontinua. La conclusión más prudente es esta: existe un narrow training interval en el que el comportamiento del modelo y sus signatures internas se reorganizan rápidamente y al mismo tiempo, lo cual justifica plantear una hipótesis sobre el mecanismo y realizar intervenciones.
\end{warningbox}

\subsection{Del cambio de comportamiento al mecanismo interno}

Si el ICL observable se fortalece de golpe en unos cuantos cientos de steps, una hipótesis natural es que dentro del modelo también se formó un nuevo circuit. Para verificarlo no basta con seguir observando la loss: hay que descomponer los parámetros en paths interpretables y comparar attention-only models de una y de dos capas, porque el primero no presenta esa phase change y el segundo sí.

\lecturefigure{slide-10-sudden-mechanism-change-hypothesis.jpg}{Hipótesis de investigación: la transición de fase en el comportamiento de ICL corresponde a la formación repentina de un algoritmo o circuito interno.}{Video oficial de Stanford Online, 00:15:12--00:15:31.}

\teachervoice{El ponente habla de una «natural hypothesis», no de una proof. La bump en el comportamiento es solo el punto de partida de la búsqueda; la explicación mecanicista solo se va sosteniendo cuando el circuit candidato aparece en el momento adecuado, su estructura puede implementar el algoritmo objetivo y la ablation elimina el comportamiento.}

\subsection{Escalera de evidencia}

Más adelante se usarán repetidamente cuatro tipos de evidencia. Para evitar confusiones, primero se ordenan por su fuerza:
\begin{table}[H]
\centering
\small
\caption{Tipos de evidencia en esta clase y las conclusiones que pueden sustentar.}
\begin{tabularx}{\textwidth}{L{0.21\textwidth}L{0.28\textwidth}X}
\toprule
Evidencia & Ejemplo & Qué puede sustentar \\
\midrule
Representational plausibility & Las matrices QK/OV pueden implementar match y copy & Muestra que el circuit tiene la capacidad, pero no demuestra que realmente se use. \\
Temporal co-occurrence & La induction signature y el ICL aparecen al mismo tiempo & Aporta evidencia correlacional; puede haber una causa común. \\
Behavioral generality & Copy, translation, pattern matching & Muestra que el mecanismo no se limita a un único toy string. \\
Causal intervention & Tras hacer ablate a la head, el ICL score retrocede & Sustenta directamente la contribución causal de ese componente a esa metric. \\
Near-complete accounting & Pocos paths explican la mayor parte de la loss del toy model & Es casi completo dentro del modelo acotado, pero no se extrapola automáticamente a LLM de gran tamaño. \\
\bottomrule
\end{tabularx}
\end{table}

\subsection{Resumen de la sección}

La phase change convierte una idea vaga, «el modelo parece aprender del context», en un objetivo de investigación concreto: explicar el cambio del ICL score dentro de una ventana estrecha. La curva de comportamiento genera la hypothesis; una conclusión mecanicista sólida requiere además álgebra, viabilidad estructural, generalización entre ejemplos y causal ablation.
